\chapter{Absorption-Line Spectroscopy: From Pixels to Physical Gas}
\label{chap:absorption}

\section{What an absorption spectrum actually measures}

An absorption spectrum is a comparison between received flux and an estimate
of the flux that would have been received without the absorber. In its most
useful forward form,
\begin{equation}
 F_{\rm pix}(\lambda)=
 \left\{
 C(\lambda)\left[(1-C_{\rm f})+
 C_{\rm f}\exp\left(-\sum_k\tau_k(\lambda)\right)\right]
 \right\}\otimes {\cal L}(\lambda)
 +B(\lambda)+\epsilon(\lambda),
 \label{eq:abs-forward}
\end{equation}
where $C$ is the source continuum and broad emission spectrum, $C_{\rm f}$ is
the covering fraction, $\mathcal{L}$ is the instrumental line-spread function
(LSF), $B$ is any residual background or scattered light, and $\epsilon$ is
noise. Each physical velocity component and each blended transition
contributes an optical-depth profile $\tau_k$. Equation~\ref{eq:abs-forward}
is a reminder that a fitted Voigt component is not directly observed. The
detector records a sampled convolution of source, absorber, instrument, and
background.

For an intervening quasar absorber the covering fraction is often taken as
unity, but this should not become reflex. Partial coverage is common in
intrinsic AGN absorption, stellar winds, compact continuum sources occulted by
small clouds, and spatially mixed emission/absorption. A nonzero flux floor in
a strong doublet can arise from partial coverage, unresolved velocity
structure, scattered light, an erroneous background, or unrelated emission.
The competing explanations must be fitted at the flux level.

The line identification itself is a hypothesis. A persuasive identification
uses multiple transitions at one redshift, expected strength ordering, absence
of plausible alternatives, and consistency with instrumental coverage.
Atomic wavelengths and oscillator strengths must come from a versioned source
such as \citet{Morton2003}; precision claims should not mix air and vacuum
wavelengths or oscillator-strength compilations silently.

\section{Coordinates, frames, and velocity conventions}

The first redshift estimate is
\begin{equation}
 1+z = \frac{\lambda_{\rm obs}}{\lambda_0}
     = \frac{\nu_0}{\nu_{\rm obs}}.
 \label{eq:redshift}
\end{equation}
For velocity structure within a system, define the convention explicitly. A
commonly used local coordinate about reference redshift $z_{\rm ref}$ is
\begin{equation}
 v=c\,\frac{z-z_{\rm ref}}{1+z_{\rm ref}}.
 \label{eq:local-velocity}
\end{equation}
At large offsets, use the exact relativistic relation or work directly in
redshift. Wavelength calibration must also specify frame: observatory,
heliocentric, barycentric, Local Standard of Rest, or a galaxy systemic frame.
Optical, radio, and relativistic velocity definitions differ at non-negligible
redshift. Combining UV, optical, and radio spectra without harmonizing
definitions can manufacture offsets comparable to the line widths being
interpreted.

The adopted systemic redshift often dominates an outflow velocity's
uncertainty, especially if it comes from a broad or asymmetric line. Narrow
stellar absorption, molecular lines, or spatially integrated nebular emission
may each define a different practical zero point.

\section{Continuum placement is part of the inference}

Dividing by a single best continuum and forgetting its uncertainty produces
overconfident line measurements. The preferred approach is to model local
continuum parameters jointly with the absorption. For an isolated line, a
low-order polynomial or spline over windows demonstrably free of absorption
may suffice. For a quasar Ly$\alpha$ forest, broad emission complexes and
correlated forest absorption require a global template, principal components,
Gaussian process, or forward model. For galaxy spectra, stellar photospheric
features and the instrumental response may be comparable to the interstellar
line.

Continuum choices should be tested by:
\begin{enumerate}
  \item changing the fit-window boundaries;
  \item increasing or decreasing continuum flexibility;
  \item injecting synthetic lines into unnormalized data;
  \item propagating continuum samples into equivalent width and column
  posteriors;
  \item checking blank regions for biased residuals and underestimated
  covariance.
\end{enumerate}
For broad, shallow features, continuum uncertainty can dominate photon noise.
A claimed broad Ly$\alpha$ absorber or AGN trough is not secure if it
disappears under plausible continuum models.

\section{Equivalent width: robust summary, incomplete physics}

For continuum-normalized flux $F_n(\lambda)$, the observed-frame equivalent
width is
\begin{equation}
 W_{\lambda,\rm obs}=\int_{\lambda_1}^{\lambda_2}
       \left[1-F_n(\lambda)\right]\,{\rm d}\lambda,
 \qquad
 W_{\lambda,\rm r}=\frac{W_{\lambda,\rm obs}}{1+z}.
 \label{eq:abs-equivalent-width}
\end{equation}
Equivalent width is conserved by a normalized LSF when all absorbed flux is
included and the continuum is correct. This makes it useful at modest
resolution. It does not make it a unique measure of column density.

In the optically thin regime, the rest equivalent width satisfies
\begin{equation}
 N\,[{\rm cm^{-2}}] =
 1.13\times10^{20}\,
 \frac{W_{\lambda,\rm r}\,[\text{\AA}]}
      {f\lambda_0^2\,[\text{\AA}^2]},
 \label{eq:abs-linear-cog}
\end{equation}
where $f$ is the oscillator strength. Once the line saturates, equivalent
width becomes sensitive mainly to Doppler width and component spread.
Eventually damping wings restore a strong column dependence. These are the
linear, saturated, and damping portions of the curve of growth
\citep{Spitzer1978,Draine2011}.

For a pixel sum with independent errors, a simple statistical variance is
\begin{equation}
 \sigma^2(W_\lambda)
 =\sum_i \left(\Delta\lambda_i\,\sigma_{F_n,i}\right)^2,
\end{equation}
but reduced spectra often have correlated noise from resampling, flat fields,
and background subtraction. Continuum uncertainty, integration-window
uncertainty, fixed-pattern noise, and blends must be added or marginalized.
An empirical distribution from line-free windows with the same width is often
more realistic than a diagonal-noise formula.

\section{Optical depth and column density}

\subsection{The apparent-optical-depth method}

For full coverage and a reliable continuum, define
\begin{equation}
 \tau_{\rm a}(v)=\ln\left[\frac{1}{F_n(v)}\right].
 \label{eq:aod-tau}
\end{equation}
The apparent column density per unit velocity is
\begin{equation}
 N_{\rm a}(v)=
 \frac{m_{\rm e}c}{\pi e^2}\frac{\tau_{\rm a}(v)}{f\lambda_0}
 =3.768\times10^{14}
 \frac{\tau_{\rm a}(v)}
      {f\lambda_0[\text{\AA}]}
 \frac{{\rm cm^{-2}}}{{\rm km\,s^{-1}}}.
 \label{eq:abs-aod-column}
\end{equation}
Integrating over velocity gives $N_{\rm a}$. The method preserves kinematic
information and enables a direct saturation test: two transitions of the same
ion should yield the same $N_{\rm a}(v)$ when resolved and unsaturated
\citep{SavageSembach1991}.

If the stronger transition yields systematically smaller $N_{\rm a}(v)$ than
the weaker one, unresolved saturation is likely. Equality is necessary but
not always sufficient: both transitions can be severely saturated and
unresolved. Negative fluxes, low signal-to-noise, and logarithmic
transformation introduce biases; simulations tailored to the data help
quantify them \citep{Fox2005}. Report lower limits rather than forcing a
precise column from black cores.

\subsection{Voigt profiles}

The intrinsic absorption cross section combines Doppler and natural
broadening. For component $j$,
\begin{equation}
 \tau_j(\nu)=N_j\frac{\pi e^2}{m_{\rm e}c}f\,
 \phi_{\rm V}(\nu;\nu_j,b_j,\Gamma),
 \label{eq:voigt-tau}
\end{equation}
where the normalized Voigt profile $\phi_{\rm V}$ is the convolution of a
Gaussian and Lorentzian, $b_j$ is the Doppler parameter, and $\Gamma$ is the
radiative damping constant. The model flux must be convolved with the
wavelength-dependent LSF before comparison with pixels.

Component decomposition is rarely unique. A high-column narrow component can
resemble a lower-column broader component after convolution. Additional weak
components can improve the likelihood without corresponding to distinct
physical clouds. Useful safeguards are:
\begin{itemize}
  \item fit all usable transitions of an ion simultaneously;
  \item tie redshifts only where a shared kinematic phase is physically
  justified;
  \item allow thermal widths to vary by ion mass rather than forcing one $b$;
  \item compare models using predictive checks, not reduced $\chi^2$ alone;
  \item report integrated columns and velocity percentiles, which are often
  more stable than the number of fitted components.
\end{itemize}

\subsection{Thermal and non-thermal broadening}

For a Gaussian component,
\begin{equation}
 b_i^2=b_{\rm nt}^2+\frac{2k_{\rm B}T}{m_i},
 \label{eq:b-decomposition}
\end{equation}
where $b_{\rm nt}$ represents unresolved turbulent, bulk, and structural
broadening. Two ions believed to be co-spatial can in principle solve for
$T$ and $b_{\rm nt}$:
\begin{equation}
 T=\frac{m_1m_2}{2k_{\rm B}(m_2-m_1)}
   \left(b_1^2-b_2^2\right).
 \label{eq:abs-two-ion-temperature}
\end{equation}
This inference fails if the ions are not co-spatial, if one profile contains
unresolved subcomponents, or if their fitted component boundaries differ.
Hydrogen plus a heavy ion gives the largest mass leverage. A broad H I line
without aligned metal absorption can constrain temperature only after blends
and continuum systematics are controlled.

\section{Saturation, doublets, and covering fraction}

The line-center optical depth for a Doppler component can be written
\begin{equation}
 \tau_0 \simeq 1.497\times10^{-15}
 \frac{N[{\rm cm^{-2}}]\,f\,\lambda_0[\text{\AA}]}
      {b[{\rm km\,s^{-1}}]}.
 \label{eq:tau-zero}
\end{equation}
This quick calculation should be made during observing design. A high
signal-to-noise spectrum does not recover information erased by saturation;
one needs weaker transitions, higher resolution, damping wings, or a
continuum edge.

For a resolved doublet with optical-depth ratio $R$ and a velocity-dependent
covering fraction,
\begin{align}
 I_{\rm weak}(v)&=1-C_{\rm f}(v)+C_{\rm f}(v)e^{-\tau(v)},\\
 I_{\rm strong}(v)&=1-C_{\rm f}(v)+C_{\rm f}(v)e^{-R\tau(v)}.
 \label{eq:partial-cover}
\end{align}
Solving these equations can separate optical depth from coverage under a
homogeneous-screen assumption. Real sources may have a distribution of
optical depth across the emitting surface, and broad-line and continuum
regions can have different coverage. Consequently, the algebraic solution is
a compact model, not a geometric image.

Doublet ratios are also distorted by blending. The O VI 1037 line, for
example, can overlap unrelated interstellar or intergalactic absorption; C IV
and Si IV doublets can overlap themselves in sufficiently broad AGN winds.
Global multiplet fitting is safer than measuring each trough independently.

\section{The curve of growth as a diagnostic}

The curve of growth predicts $W_\lambda(N,b)$ for one component. Fitting
several transitions with different $f\lambda$ can constrain $N$ and $b$ even
when profiles are not resolved. The method is valuable for sparse data but
compresses away velocity information. Multi-component structure produces an
effective $b$ whose physical meaning is limited.

A practical curve-of-growth analysis should display the data and posterior in
$W/\lambda$ versus $Nf\lambda$ space, list the transitions excluded for
blending, and test whether a single $(N,b)$ fits both weak and strong lines.
If only saturated transitions are available, the posterior often has a long
column-density tail. Quoting symmetric errors in that situation is
misleading. The output may properly be a lower bound.

\section{Blends and line confusion}

Line confusion increases toward short rest wavelength, high source redshift,
and low resolution. A candidate feature should be tested against a line list
at every known absorber redshift, Galactic ISM and molecular features,
geocoronal emission, detector artifacts, and telluric absorption for
ground-based data. The most reliable workflow is simultaneous:
\begin{enumerate}
  \item construct a master list of expected transitions for all systems;
  \item map every transition into observed wavelength;
  \item fit overlapping optical depths in one flux model;
  \item require shared parameters only within physically motivated groups;
  \item inspect residuals in detector coordinates as well as wavelength.
\end{enumerate}

Airglow can be time dependent; separating orbital night may reduce it at the
cost of exposure. Telluric absorption is multiplicative and depends on
airmass, humidity, and time. Sky emission is additive and has a spatial
profile. Treating all three with a generic mask discards information and can
leave biased wings. A forward model or physically chosen exclusion is
preferable.

\section{Instrumental resolution and sampling}

Resolving power $R=\lambda/\Delta\lambda$ corresponds approximately to
\begin{equation}
 \Delta v_{\rm FWHM}\simeq \frac{c}{R}.
 \label{eq:r-to-v}
\end{equation}
The equation is a planning shorthand. Real LSFs can be non-Gaussian,
wavelength dependent, and altered by source extent, slit illumination, focus,
or detector location. A model using a Gaussian with the nominal $R$ can bias
narrow-line columns and component widths. Use the tabulated or empirical LSF
for the actual setting.

Nyquist sampling is not the same as spectral resolution. Rebinning can improve
visual signal-to-noise per bin but does not create information and generates
covariance. Keep native or minimally resampled pixels for the likelihood;
rebin for presentation only, or carry the covariance matrix. Coaddition is
also a model: exposures can differ in wavelength zero point, LSF, sensitivity,
fixed-pattern structure, and background. Fitting exposures jointly is often
safer than coadding first.

\section{From ionic columns to a physical phase}

\subsection{Photoionization models}

Define the ionization parameter
\begin{equation}
 U=\frac{\Phi_{\rm H}}{n_{\rm H}c}
  =\frac{n_\gamma(E>13.6\,{\rm eV})}{n_{\rm H}}.
 \label{eq:abs-ionization-parameter}
\end{equation}
Given an incident spectrum, elemental abundances, geometry, stopping column,
and equilibrium assumptions, a code such as Cloudy predicts ionic columns
\citep{Ferland2017}. Ratios of successive stages of one element can constrain
$U$ while reducing abundance sensitivity. Combining H I with metal ions then
constrains metallicity. In practice, multiphase structure often prevents one
$U$ from fitting every ion.

The radiation field is part of the model, not a universal constant.
Metagalactic UV-background prescriptions differ in amplitude and hardness
\citep{HaardtMadau2012}; nearby galaxies or an AGN can dominate locally.
Because $U$ constrains a photon-density to gas-density ratio, changing the
field changes the inferred $n_{\rm H}$ and path length
\begin{equation}
 L=\frac{N_{\rm H}}{n_{\rm H}}.
\end{equation}
Path lengths that exceed the halo or imply implausible pressure are powerful
model checks.

A reproducible photoionization analysis reports the code version, atomic
database, incident spectrum and normalization, solar abundance reference,
relative-abundance priors, dust/depletion treatment, geometry, density law,
stopping rule, and whether upper limits enter the likelihood. Publishing the
model grid or scripts is more useful than publishing a corner plot alone.

\subsection{Collisional and non-equilibrium models}

In collisional ionization equilibrium (CIE), ion fractions depend mainly on
temperature for a specified abundance set. This makes ion ratios appear to be
thermometers. The assumptions are restrictive: the gas must have had time to
ionize and recombine, photoionization must be negligible, and electrons must
be close to a Maxwellian distribution. Rapidly cooling gas can remain
overionized, while a shock can remain underionized. Non-equilibrium cooling
calculations show that the same ionic ratio may correspond to different
thermal histories \citep{GnatSternberg2007}.

The cooling time and recombination time should be compared:
\begin{equation}
 t_{\rm rec}(X^i)=\frac{1}{\alpha_{i}(T)n_{\rm e}}.
\end{equation}
If $t_{\rm cool}\lesssim t_{\rm rec}$, equilibrium is doubtful. Turbulent
mixing and conduction create distributions of temperature rather than one
$T$. In that case, integrated interface columns and kinematic alignment may
be more meaningful than a single-temperature fit.

\section{Abundances, depletion, and ionization corrections}

An elemental abundance inferred from ions is
\begin{equation}
 [{\rm X/H}]=
 \log\left(\frac{N_{X^i}}{N_{\rm H\,I}}\right)
 +\log\left(\frac{f_{\rm H\,I}}{f_{X^i}}\right)
 -\log({\rm X/H})_\odot.
 \label{eq:abs-ionic-abundance}
\end{equation}
Every term needs a convention. The solar scale changes at the
$0.01$--$0.1$ dex level for important elements \citep{Asplund2021}. Dust
removes refractory elements from the gas. Nucleosynthetic patterns can depart
from solar, especially at low metallicity or after delayed Type Ia and AGB
enrichment. Charge exchange can couple some ions to H I. Treating every
difference as ionization produces a deceptively precise but wrong
metallicity.

Useful abundance combinations include:
\begin{itemize}
  \item O I/H I for predominantly neutral gas, with charge-exchange caveats;
  \item S II or Zn II as comparatively weakly depleted tracers where usable;
  \item Fe/alpha-element ratios for depletion plus enrichment history;
  \item N/O for secondary versus primary nitrogen, provided ionization phases
  match.
\end{itemize}
Upper limits are scientifically active. A non-detection of a low ion can rule
out a high-density solution; a non-detection of N V can disfavor one hard
spectrum. Censoring limits from a likelihood because they are not detections
throws away this leverage.

\section{Population measurements}

For an absorber survey, completeness is a function
$C(W,z,\lambda,b,\mathrm{S/N},\ldots)$, not one threshold. The incidence is
often written
\begin{equation}
 \frac{{\rm d}{\cal N}}{{\rm d}z}
 =\frac{1}{\Delta z_{\rm eff}}
  \sum_j \frac{1}{C_j},
 \qquad
 \Delta z_{\rm eff}=\int g(z)\,{\rm d}z,
\end{equation}
where $g(z)$ accounts for coverage blocked by detector gaps, foreground
features, proximity exclusions, and sensitivity. For cosmological comparison,
one can use absorption distance
\begin{equation}
 {\rm d}X=\frac{(1+z)^2H_0}{H(z)}\,{\rm d}z.
\end{equation}
Blind injection and recovery through the actual detection pipeline is the
most direct completeness calibration.

Association with galaxies introduces another selection function: magnitude
limit, angular coverage, redshift success, fiber collision, and group-finding
algorithm. Calling an absorber isolated when the galaxy survey is incomplete
is not a physical result. The low-redshift COS survey literature provides
useful examples of documenting redshift path and line-finding methodology
\citep{Danforth2016}.

\section{A reproducible analysis recipe}

\begin{enumerate}
  \item Preserve calibrated exposures and quality arrays; do not begin from a
  screenshot or display-only coadd.
  \item Verify wavelength frames and align exposures using narrow,
  well-identified features without erasing real shifts.
  \item Define a transition list with versioned atomic data.
  \item Fit continuum, background, absorption, and the correct LSF jointly
  where practical.
  \item Measure equivalent widths and apparent columns as minimally
  model-dependent summaries.
  \item Test saturation with weaker lines, doublets, higher Lyman lines, or
  Lyman-edge and damping information.
  \item Fit physically motivated multiplets, but separate robust integrated
  quantities from non-unique component structure.
  \item Compare photoionized, collisionally ionized, and non-equilibrium
  hypotheses; use posterior predictive spectra.
  \item Propagate atomic, continuum, calibration, radiation-field, abundance,
  depletion, and geometry systematics into derived mass or energetics.
  \item Release normalized and unnormalized spectra, masks, LSFs, code, model
  grids, posterior samples, and machine-readable measurements.
\end{enumerate}

\section{Failure modes worth remembering}

\begin{description}
  \item[Black core, precise column:] impossible without weaker transitions or
  damping and edge information; the correct result is often a limit.
  \item[Same velocity, same phase:] not guaranteed in a projected,
  multi-component medium.
  \item[Large $b$, high temperature:] only an upper limit if non-thermal and
  unresolved broadening are possible.
  \item[Doublet ratio looks right:] both members may be unresolved and
  saturated, or one may be blended.
  \item[One-zone model fits most ions:] the excluded ion may be the clearest
  evidence for a second phase.
  \item[High formal significance:] unmodeled continuum flexibility, line
  searching across many resolution elements, and calibration features change
  the false-positive rate.
  \item[Precise mass from one sightline:] geometry and covering fraction can
  dominate the error even when the column is secure.
\end{description}

Absorption spectroscopy is powerful because it detects very little material.
That same sensitivity makes it vulnerable to hidden saturation, unrelated
blends, and model-dependent ionization corrections. The defensible strategy
is layered: preserve the pixel-level evidence, extract model-light
measurements, and add physical assumptions one at a time.
